% Generated draft: requires mathematical and source review.
\documentclass[11pt]{article}
\usepackage[margin=25mm]{geometry}
\usepackage{fontspec}
\setmainfont{Latin Modern Roman}
\usepackage{amsmath,amssymb,mathtools,unicode-math}
\setmathfont{Latin Modern Math}
\usepackage{enumitem,needspace,xurl,hyperref}
\setlength{\parindent}{0pt}
\setlength{\parskip}{5pt}
\setlength{\emergencystretch}{3em}
\newcommand{\N}{\mathbb{N}}
\newcommand{\Z}{\mathbb{Z}}
\newcommand{\Q}{\mathbb{Q}}
\newcommand{\R}{\mathbb{R}}
\newcommand{\C}{\mathbb{C}}
\newcommand{\F}{\mathbb{F}}
\newcommand{\A}{\mathbb{A}}
\newcommand{\PP}{\mathbb{P}}
\newcommand{\E}{\mathbb{E}}
\newcommand{\eps}{\varepsilon}
\newcommand{\dd}{\,\mathrm{d}}
\begin{document}
\section*{UnsolvedMath ID: 701116}
\textbf{Source problem code:} catalogue entry
\subsection{Definitions and mathematical statement}
Let \(p\) be a prime and let \(G\) be a finite \(p\)-group, meaning that \(|G|=p^n\) for some integer \(n\geq 0\). For a finite group \(H\), define its exponent by
\[
\exp(H)=\min\{m\in\Z_{>0}:h^m=1\text{ for every }h\in H\}.
\]
Let \(A,B\leq G\) satisfy
\[
G=AB=\{ab:a\in A,\ b\in B\}.
\]
No normality assumption on either subgroup, or restriction on their intersection, is supplied.

\begin{enumerate}
\item[(a)] Does there exist a function
\[
F:\Z_{>0}\times\Z_{>0}\longrightarrow\Z_{>0}
\]
such that, for every prime \(p\), every finite \(p\)-group \(G\), and all subgroups \(A,B\leq G\) with \(G=AB\),
\[
\exp(G)\leq F\bigl(\exp(A),\exp(B)\bigr)?
\]
Thus the bound may depend only on the two subgroup exponents, with no further dependence on the groups or their orders.
\item[(b)] Is \(\exp(G)\) bounded when \(\exp(A)=\exp(B)=p\)? With dependence on the prime allowed, the precise question is whether
\[
\begin{gathered}
\forall\,p\text{ prime}\ \exists\,C_p\in\Z_{>0}\ \forall\,G\text{ a finite }p\text{-group}\ \forall\,A,B\leq G,\\
\bigl(G=AB\ \text{and}\ \exp(A)=\exp(B)=p\bigr)
\ \Longrightarrow\ \exp(G)\leq C_p.
\end{gathered}
\]
The source does not explicitly specify the permitted dependence of the bound in part (b). The displayed formulation allows dependence on \(p\); a bound independent of \(p\) would instead require a single constant \(C\) preceding the universal quantifier over primes. This qualification remains unresolved from the supplied record.
\end{enumerate}
\subsection{Short English statement}
For a finite \(p\)-group expressed as a product \(G=AB\) of two subgroups, ask whether the exponent of \(G\) is bounded using only the exponents of \(A\) and \(B\). Also ask whether such a bound exists when both subgroup exponents equal \(p\), with the permitted dependence on \(p\) requiring clarification.
\subsection{Sources}
\begin{itemize}
\item Local UnsolvedMath registry: supplied complete record with permanent ID \(701116\), exact source problem code \texttt{catalogue entry}, title \emph{Kourovka Notebook Problem 15.28}. The record was created and updated on 2026-10-16; its upstream revision is \texttt{d371ee917ac9647fbb06cd6fc529deec6f3dda82}.
\item Supplied source attribution: \emph{The Kourovka Notebook (Unsolved Problems in Group Theory)}, Problem 15.28, proposed by Ya. G. Berkovich in 2002, in the 15th issue (2002). The record also identifies the 21st issue edition (2026), \texttt{arXiv:1401.0300v48}. The collection description names E. I. Khukhro and V. D. Mazurov as editors.
\item The registry labels this problem \emph{open}. This is a supplied status label, not an independently verified assessment. The cited editions and their contents have not been independently checked for this draft.
\end{itemize}
\end{document}
